\section{Faithful tensor tests in relative Chow}
\label{sec:tensors}

The formal detector and the Chow-vanishing argument use different
coefficient algebras.  For the latter we must retain all Chow
correspondences, including those not generated by the pairing.  We
construct such a tensor model in this section.  Its coefficient algebra
is generally large; the point is its faithfulness, not a presentation of
its generators.

Let
\[
 \pi\colon C\longrightarrow S
\]
be a smooth projective family of connected curves of genus \(b\), with a
section, over an irreducible smooth quasi-projective complex variety.
Fix a rational divisor class \(e_C\in\CH^1(C)\) of relative degree one.
The class \(e_C\) need not be the class of the section.  Set
\begin{equation}
 \label{eq:curve-projector}
 a_C=\pi_*(e_C^2),\qquad
 e'_C=e_C-\frac12\pi^*a_C,\qquad
 \eta=[\Delta_{C/S}]-e'_{C,1}-e'_{C,2}.
\end{equation}
Here and below products of curve families are over \(S\).

For an integer \(a>0\), we may work modulo classes that vanish after
restriction to \(S\setminus Z\), for some closed subset \(Z\subset S\)
of codimension at least \(a\).  This convention does not mean taking
numerical equivalence, nor does it mean testing individual geometric
fibers.

\begin{proposition}[Faithful Chow tensor test]
 \label{prop:tensor-test}
In the setup above, either in relative rational Chow or in the indicated
quotient by base support, there is a supercommutative rational algebra
\(\mathfrak B\) with the following properties.
\begin{enumerate}[label=\textup{(\roman*)}]
\item
The curve correspondence algebra has a faithful tensor model on
\[
 \mathfrak B\otimes
 \bigl(\Q 1\oplus W_b\oplus\Q e\bigr),
\]
where \(1,e\) are even and \(W_b\) is an odd symplectic space of dimension
\(2b\).  The classes \(1_C,e'_C\) correspond to \(1,e\), and \(\eta\)
projects onto \(W_b\).  Integration gives the hyperbolic pairing
\(\langle1,e\rangle=1\), \(\langle1,1\rangle=\langle e,e\rangle=0\)
on the two even summands and the symplectic pairing on \(W_b\).
All actual Chow correspondences between curve powers, not only pairing
diagrams, are included faithfully in this model.
\item
The projected small diagonal is an alternating three-form
\[
 \varphi\colon W_b^3\longrightarrow\mathfrak B_{\mathrm{odd}}.
\]
Cup product, integration, tensor products, permutations, and contractions
are computed in the model with the Koszul signs.
\item
Let \(U\) be a finite-dimensional ordinary rational color space.  A pure
colored tensor in
\begin{equation}
 \label{eq:colored-chow-space}
 \left(
   \eta^{\otimes k}\CH^*(C^k)
   \otimes U^{\otimes k}\otimes\mathrm{sgn}_k
 \right)_{\mathfrak S_k}
\end{equation}
is zero if and only if all its evaluations using the same ordinary
linear map \(z\colon U\to W_b\) in every slot are zero in \(\mathfrak B\).
Finitely many rational maps \(z\) suffice for each fixed \(U,k\).
The assertion holds in the base-support quotient as well.
\end{enumerate}
In particular, a product of projected cubic entries is tested by
products of
\(\varphi(z(i),z(j),z(k))\), with the fixed Koszul sign coming from the
ordering of the factors.  A vanishing established by these tests is a
vanishing in Chow, with the stated base-support bound.
\end{proposition}

\subsection{Correspondences and the curve projector}

We use the rational correspondence category whose objects are smooth
projective \(S\)-families and whose morphisms are
\[
 \Hom(P,Q)=\CH^*(P\times_S Q),
\]
with all codimensions allowed.  Composition is pullback to the triple
product, intersection, and pushforward.  Tensor product is relative
product.  Equivalently, this is the usual correspondence calculus with
Tate twists forgotten.  We take finite direct sums and split
idempotents.  The diagonal identifies every generating object with its
dual; the evaluation for the curve is its intersection pairing followed
by \(\pi_*\).

The morphisms that vanish away from a closed base subset of codimension
at least \(a\) form a tensor ideal.  Indeed, a finite union of such
subsets has the same codimension bound, and restriction to its
complement commutes with every operation just described.  We can
therefore first quotient by this ideal and then split idempotents.
Morphisms between the original objects are unchanged by the latter
completion.  Denote either the unquotiented category or this completed
quotient by \(\mathcal C\).
Its tensor unit is the base \(S\).  The endomorphisms of that unit are
\(\CH^*(S)\), or their base-support quotient; positive-degree base
classes remain part of the correspondence calculus unless the chosen
quotient kills them.

The projection formula gives
\begin{equation}
 \label{eq:hyperbolic-normalization}
 \pi_*1_C=0,\qquad
 \pi_*e'_C=1,\qquad
 \pi_*((e'_C)^2)=a_C-a_C=0.
\end{equation}
Thus \(1_C,e'_C\) span a nondegenerate hyperbolic plane in the curve
object.  The two corresponding projectors sum to
\(e'_{C,1}+e'_{C,2}\), so their orthogonal complement has the
self-transpose projector \(\eta\).  Write this complementary object as
\(H_C\).

The categorical trace of the curve identity is
\(\pi_*c_1(T_{C/S})=2-2b\), by diagonal self-intersection.  Each of the
two unit summands has trace one.  Consequently
\begin{equation}
 \label{eq:odd-curve-dimension}
 \dim_{\mathcal C}H_C=-2b.
\end{equation}
The remaining relation needed to model \(H_C\) by an odd symplectic
space holds already in relative Chow.

\begin{lemma}[Relative pairing relation]
 \label{lem:relative-pfaffian}
On \(C^{2b+2}\), one has
\begin{equation}
 \label{eq:relative-pfaffian}
 \sum_{\mathcal M}
   \prod_{\{i,j\}\in\mathcal M}\eta_{ij}=0,
\end{equation}
where \(\mathcal M\) ranges over the perfect matchings of
\(\{1,\ldots,2b+2\}\).
\end{lemma}

\begin{proof}
Let \(A\to S\) be the relative Jacobian. The section of \(C\) permits a
universal degree-zero line bundle on \(C\times_S A\), normalized along
that section \cite[Section~9.4, Proposition~4; Section~8.2,
Proposition~4]{BLR}. We use the Poincar\'e bundle and biduality for
abelian schemes in \cite[Section~8.4, Theorem~5]{BLR}.  If \(\ell_1,\ell_2\) are its two first Chern classes on
\(C\times_S A\times_S A\), put
\[
 P=(\pi_{A\times A})_*(\ell_1\ell_2)\in\CH^1(A\times_S A),
 \qquad u=\frac12\Delta_A^*P\in\CH^1(A).
\]
Base twists of the universal line do not affect \(P\), because its
relative degree is zero.  The class \(P\) is symmetric and biadditive.
In particular
\begin{equation}
 \label{eq:jacobian-eigenclass}
 [m]^*u=m^2u.
\end{equation}
We first prove \(u^{b+1}=0\) in relative rational Chow.

The Fourier method originates in the work of Mukai and Beauville
\cite[Theorem~2.2]{Mukai1981}, \cite[Section~2, Proposition~3]{Beauville1983}.
We give the argument with its relative base factor.
For the normalized Poincare line on \(A\times_S\widehat A\), write its
first Chern class as \(c\), and define on total Chow
\[
 \mathcal F(\alpha)=
 p_{\widehat A*}\bigl(p_A^*\alpha\,\exp(c)\bigr).
\]
All exponential expressions are finite in each Chow ring.  Convolution
with the analogous transform for \(\widehat A\) has kernel
\[
 (+)^*K,\qquad
 K=p_{A*}\exp(c),
\]
on \(A\times_S A\).  This follows by the biadditivity of the Poincare
class and flat base change.

The class \(K\) is supported on the zero section \(0_A\).  To see this,
apply Grothendieck--Riemann--Roch \cite[Section~7]{BorelSerre} to the Poincare line under
\(p_A\).  The relative tangent bundle of an abelian scheme is pulled
back from the base, so its Todd factor is an invertible base class.
The derived pushforward of the Poincare line vanishes away from
\(0_A(S)\): on a geometric fiber, a nontrivial degree-zero line on an
abelian variety has zero cohomology, and perfect direct image and derived base change give the assertion
\cite[Tags~0B91 and~068V]{Stacks}.  The fiber statement can
also be seen from the flat unitary-character description of
\(\Pic^0\): a nontrivial character has no harmonic Dolbeault forms on a
flat torus.

Chow localization now gives \(K=0_{A*}\delta\) for
\(\delta\in\CH^*(S)\).  This coefficient is unique: if
\(\pi_A\colon A\to S\), then \(\pi_{A*}0_{A*}=\id\), so
\(\delta=\pi_{A*}K\).  Its degree-zero component is nonzero.  It is
the fiber integral
\[
 \int_{A_s\times\widehat A_s}\frac{c^{2b}}{(2b)!},
\]
which is nonzero because the Poincare class is the full-rank evaluation
pairing between the two first cohomology spaces.  This computation
determines only a scalar, not a Chow class by its realization.
The positive-degree ideal of \(\CH^*(S)\) is nilpotent, so \(\delta\)
is invertible.  Flat base change along addition identifies
\((+)^*0_{A*}\delta\) with the graph of inversion carrying the
pullback of \(\delta\).  Convolution therefore gives inversion on
\(A\), times this invertible base class.  In particular
\(\mathcal F\) is injective.

The same biadditivity and projection formula show
\[
 \mathcal F\circ[m]_*
   =[m]^*\circ\mathcal F.
\]
Furthermore its component
\[
 \mathcal F_j(\alpha)=
 p_{\widehat A*}\left(p_A^*\alpha\,\frac{c^j}{j!}\right)
\]
has \([m]^*\)-eigenvalue \(m^j\), where \(j\ge0\).
On the other hand, \([m]\) is finite flat of degree \(m^{2b}\), and
\eqref{eq:jacobian-eigenclass} gives
\[
 [m]_*u^{b+1}=m^{-2}u^{b+1}.
\]
For a fixed integer \(m>1\), the number \(m^{-2}\) is distinct from
every eigenvalue \(m^j\) occurring in the finite sum
\(\mathcal F(u^{b+1})\).  Applying the product of the corresponding
linear polynomials in \([m]^*\) shows that this transform is zero.
Injectivity proves the claimed vanishing.

Choose \(N>0\) clearing the denominators of \(e_C\).  The maps
\[
 a_N\colon C\longrightarrow A,\qquad
 x\longmapsto[Nx-Ne_C]
\]
are defined relatively.  On an additional curve coordinate, the
pullback of the universal degree-zero line has class
\(N[\Delta]-Ne_C\), up to a twist from the parameter curve.
Integrating the product for two parameter positions gives
\begin{equation}
 \label{eq:abel-jacobi-cross-term}
 (a_N\times a_N)^*P
 =N^2\bigl([\Delta]-e_{C,1}-e_{C,2}+\pi^*a_C\bigr)
 =N^2\eta.
\end{equation}
The possible base twists contribute zero because the other line has
relative degree zero.

Pull back \(u^{b+1}=0\) along
\[
 (x_1,\ldots,x_{2b+2})
 \longmapsto\sum_i t_i a_N(x_i),
\]
where the \(t_i\) are independent integers.  Biadditivity gives a
polynomial identity in these integers.  The coefficient of
\(t_1\cdots t_{2b+2}\) is
\[
 (b+1)!\,N^{2b+2}
 \sum_{\mathcal M}\prod_{\{i,j\}\in\mathcal M}\eta_{ij}.
\]
Indeed, the self terms have a squared multiplier and cannot
contribute to this coefficient.  Rational polynomial interpolation
and division by the displayed nonzero scalar prove
\eqref{eq:relative-pfaffian}.
\end{proof}

\subsection{Retaining arbitrary Chow correspondences}

Let \(\mathcal R_b=\operatorname{Rep}^{\mathrm{odd}}(\Sp_{2b})\).
This is the usual semisimple rational representation category of the
split symplectic group, with parity prescribed by its central
involution.  Its standard representation \(W_b\) is odd.  The ordinary
alternating symplectic form is therefore a symmetric pairing in the
super sense; its closed loop has value \(-2b\).

The first and second fundamental theorems for symplectic tensors
present this category, after finite direct sums and idempotent
completion, by pairing diagrams with that loop value and the
sum-of-matchings relation on \(2b+2\) slots.  In this formulation
crossings on the standard odd space are minus the ordinary swap.
See \cite[Theorems~2.2 and~2.13]{RW15}, or
\cite[Theorems~4.6, 4.8, and~5.9]{LZ15}.
The diagram spaces and relations are rational, so the complex
formulation descends by faithful scalar extension to \(\Q\).
Complete reducibility, absolute simplicity, and tensor generation over
\(\Q\) follow from \cite[Theorems~4.14, 22.2, and~22.42, and
Paragraph~22.4]{Milne2022}. In particular, the irreducibles occur as
rational direct summands of tensor powers of the standard representation.
Equations
\eqref{eq:odd-curve-dimension} and \eqref{eq:relative-pfaffian} thus
give a symmetric tensor functor
\begin{equation}
 \label{eq:symplectic-chow-functor}
 F\colon\mathcal R_b\longrightarrow\mathcal C,
 \qquad F(W_b)=H_C.
\end{equation}
For \(b=0\), \eqref{eq:relative-pfaffian} says \(\eta=0\), which gives
the same conclusion with \(W_0=0\).

It is essential not to assume that \eqref{eq:symplectic-chow-functor}
is full.  The target contains the small diagonal and many other
correspondences.  We record them by enlarging the coefficient algebra.
The following reconstruction is a concrete split semisimple form of
O'Sullivan's tensor realization
\cite[Lemma~2.3.1\textup{(i)}]{OSullivan2011}.  It applies to a
symmetric tensor functor with this semisimple source and an arbitrary
rational tensor target; the endomorphism ring of the target unit need
not be a field.  We give its coefficient construction explicitly.

\begin{lemma}[Coefficient algebra]
 \label{lem:coefficient-algebra}
Let \(F\) be the functor in \eqref{eq:symplectic-chow-functor}, and put
\(H(X)=\Hom_{\mathcal C}(1,F(X))\).  There is a supercommutative
\(\Sp_{2b}\)-algebra \(\mathfrak B\) such that, naturally in \(X\),
\begin{equation}
 \label{eq:coefficient-invariants}
 H(X)\simeq(\mathfrak B\otimes X)^{\Sp_{2b}}.
\end{equation}
Under these identifications, tensor products of target morphisms
become multiplication of coefficients.  Arbitrary morphisms
\(F(X)\to F(Y)\) are represented faithfully by equivariant
\(\mathfrak B\)-linear maps
\(\mathfrak B\otimes X\to\mathfrak B\otimes Y\), and their composition
is the usual tensor contraction.
\end{lemma}

\begin{proof}
Choose one representative \(X_\lambda\) of each irreducible rational
symplectic representation and set
\begin{equation}
 \label{eq:coefficient-algebra-definition}
 \mathfrak B=
 \bigoplus_\lambda H(X_\lambda)\otimes X_\lambda^\vee.
\end{equation}
The multiplicity spaces \(H(X_\lambda)\) are ordinary vector spaces;
the parity on each summand comes from \(X_\lambda^\vee\).
Semisimplicity gives
\[
 H(X)\simeq
 \bigoplus_\lambda
 H(X_\lambda)\otimes\Hom_{\Sp_{2b}}(X_\lambda,X),
\]
which is exactly \eqref{eq:coefficient-invariants}.
Only finitely many summands enter for any fixed \(X\).

An element \(s\in H(X)\) corresponds equivalently to an equivariant
map \(i_s\colon X^\vee\to\mathfrak B\).  For
\(s\in H(X)\) and \(t\in H(Y)\), the target tensor product gives
\[
 s\otimes t\in H(X\otimes Y).
\]
Define multiplication on \(\mathfrak B\) by requiring the product of
\(i_s\) and \(i_t\) to be the map corresponding to this element,
using the natural super identification
\(X^\vee\otimes Y^\vee\simeq(X\otimes Y)^\vee\).
Explicitly, for homogeneous dual vectors \(f,g\) this identification
sends \(f\otimes g\) to
\[
 x\otimes y\longmapsto(-1)^{|g||x|}f(x)g(y).
\]
On the irreducible summands of
\eqref{eq:coefficient-algebra-definition} this specifies every
product; naturality and semisimplicity give the same rule for all
\(X,Y\).  Associativity, the unit, and supercommutativity follow from
the associativity, unit, and symmetry constraints for the target
tensor product.  Thus \eqref{eq:coefficient-invariants} respects
tensor products, including their Koszul signs.

Finally rigidity identifies
\[
 \Hom_{\mathcal C}(F(X),F(Y))
 \simeq H(X^\vee\otimes Y).
\]
Apply \eqref{eq:coefficient-invariants} and the usual tensor--Hom
identification for finite free supermodules:
\[
 H(X^\vee\otimes Y)
 \simeq(\mathfrak B\otimes X^\vee\otimes Y)^{\Sp_{2b}}
 \simeq
 \Hom^{\mathrm{even}}_{\mathfrak B,\Sp_{2b}}
   (\mathfrak B\otimes X,\mathfrak B\otimes Y).
\]
This gives the asserted faithful representation of every target
morphism.  Evaluation and
coevaluation are images of the corresponding representation
morphisms, so composition is precisely coefficient multiplication
and contraction.  Faithfulness here is the injectivity of an actual
vector-space isomorphism, not faithfulness of a cohomological
realization.
\end{proof}

Apply the lemma to
\[
 V_b=\Q1\oplus W_b\oplus\Q e,
 \qquad F(V_b)\simeq C.
\]
In particular, a pure Chow class on \(C^k\) is a morphism
\(1\to H_C^{\otimes k}\), so it is identified with an invariant in
\(\mathfrak B\otimes W_b^{\otimes k}\).  This identification retains
the class itself, including any component invisible to cohomology.
The small diagonal and the structure map supply multiplication,
unit, and integration.  Their correspondence identities give a
supercommutative Frobenius algebra on \(\mathfrak B\otimes V_b\).
The pairing is the one fixed by
\eqref{eq:hyperbolic-normalization}.  In particular
\begin{equation}
 \label{eq:actual-cubic-form}
 \varphi(w_1,w_2,w_3)=\int_C w_1w_2w_3
 \quad (w_i\in W_b)
\end{equation}
is the tensor of the small diagonal after projection by \(\eta\) in
all three slots.  Its coefficients are odd, and it is alternating
in the ordinary underlying vectors \(w_i\). For a pointed curve over a
field, this projection is the modified diagonal of Gross and Schoen
\cite[Equation~(0.1) and Sections~2--3]{GrossSchoen1995}. The normalization
here also incorporates the base corrections required for a family.

\begin{lemma}[Centering at a section]
 \label{lem:point-centered-algebra}
Work modulo positive-codimension base support, and suppose \(e_C\)
is the class of a section.  Then \(e'_C=e_C\) in the quotient.
In the model of Proposition~\ref{prop:tensor-test}, the functional
\[
 \epsilon(x)=\int_C e_Cx
\]
is multiplicative and satisfies
\[
 \epsilon(1)=1,\qquad
 \epsilon(e)=0,\qquad
 \epsilon(W_b)=0.
\]
Consequently the component of a product \(w_1w_2\), with
\(w_1,w_2\in W_b\), in the two even summands is
\(\langle w_1,w_2\rangle e\).
\end{lemma}

\begin{proof}
Every positive-degree base class vanishes after removing a suitable
positive-codimension subset.  In particular \(a_C\) vanishes, so
\(e'_C=e_C\).  Pairing with \(e_C\) is pullback along the section, by
the projection formula, and this pullback is multiplicative.
The displayed values follow from the hyperbolic splitting and the
orthogonality of \(H_C\).

The coefficient of \(1\) in \(w_1w_2\) is its pairing with \(e\),
namely \(\epsilon(w_1w_2)=0\).  The coefficient of \(e\) is its
pairing with \(1\), namely \(\langle w_1,w_2\rangle\).  These are
identities of actual correspondences, hence also identities in the
faithful coefficient model.
\end{proof}

\subsection{Colored operations and finite polarization}

For an ordinary rational color space \(U\), define
\[
 \mathcal A_k(U)=
 \left(
  \CH^*(C^k)\otimes U^{\otimes k}\otimes\mathrm{sgn}_k
 \right)_{\mathfrak S_k}.
\]
The product is pullback to disjoint curve slots followed by the
coinvariant projection.  With parity \(k\bmod2\) in arity \(k\), the
direct sum of these spaces is supercommutative: interchanging blocks
of \(k\) and \(\ell\) slots has sign \((-1)^{k\ell}\).
The same definition applies to a collection of configuration spaces
with equivariant forgetting maps, using their forgetting pullbacks
for the product.  We will use that version for weighted pointed
curves.

A symmetric three-slot class gives an odd alternating entry
\(f(i,j,k)\), multilinear in its colors.  For the pure small diagonal,
its scalar-color evaluation is
\(\varphi(z(i),z(j),z(k))\).  We use a fixed ordering of tensor factors
throughout; reordering several entries introduces exactly their
Koszul signs.

\begin{lemma}[Finite color polarization]
 \label{lem:color-polarization}
In the pure subspace
\(\eta^{\otimes k}\mathcal A_k(U)\), evaluations using the same scalar
map \(z\colon U\to W_b\) in every slot are jointly injective.
For fixed \(k,U\), a finite collection of rational maps is jointly
injective.  These statements hold after quotienting by base support
of any prescribed positive codimension.
\end{lemma}

\begin{proof}
By Lemma~\ref{lem:coefficient-algebra}, before taking color
coinvariants a pure class is an invariant tensor with coefficients
in \(\mathfrak B\) and \(k\) odd standard factors.  The super
permutation of those factors is the ordinary permutation multiplied
by \(\mathrm{sgn}_k\).  It cancels the explicit sign representation
in the definition of \(\mathcal A_k(U)\).
Writing \(W_{b,\mathrm{ord}}\) for the underlying ordinary vector
space, the invariant-tensor inclusion therefore induces an injection
\[
 \eta^{\otimes k}\mathcal A_k(U)
 \lhook\joinrel\longrightarrow
 \mathfrak B\otimes
 \Sym^k(W_{b,\mathrm{ord}}\otimes U).
\]
The symmetric group acts trivially on the single coefficient factor
\(\mathfrak B\) in this formula.

Using the nondegenerate symplectic pairing to identify a standard
factor with its dual, the colored tensor is therefore an ordinary
homogeneous polynomial tensor of degree \(k\) on
\(\Hom_\Q(U,W_b)\), with coefficients in \(\mathfrak B\).
The invariant-tensor identification is injective, and taking finite
symmetric-group coinvariants is exact over \(\Q\).  Hence no
additional kernel was introduced in obtaining this polynomial
tensor.

Evaluation at \(z\) is precisely polynomial evaluation.  A
homogeneous polynomial with coefficients in any rational vector
space is determined by its values on rational points.  More
explicitly, after choosing bases, interpolation on a finite rational
grid recovers every degree-\(k\) coefficient.  The interpolation
matrix has rational entries and a rational left inverse.  It
therefore works without requiring the coefficients to be reduced,
even, or finite-dimensional.  This proves both injectivity
assertions.

The proof can be carried out from the start in the quotient
correspondence category.  Alternatively, use the same finite
interpolation: take the union of the finitely many closed base
supports occurring in the evaluations and coefficient
identifications.  Its codimension is at least the prescribed bound.
Restriction to its complement makes every evaluation zero, and the
interpolation identity makes the original colored class zero
there.  This explains why a separate exceptional support for each
test does not weaken the conclusion.
\end{proof}

\begin{proof}[Proof of Proposition~\ref{prop:tensor-test}]
The hyperbolic decomposition, Lemma~\ref{lem:relative-pfaffian}, and
the symplectic diagram presentation construct the functor \(F\).
Lemma~\ref{lem:coefficient-algebra} supplies the faithful model for
all curve correspondences and therefore for the Frobenius algebra.
Equation~\eqref{eq:actual-cubic-form} identifies the projected small
diagonal.  Lemma~\ref{lem:color-polarization} supplies the
same-map test and its base-support assertion.  Tensor products
represent products of disjoint-slot operations, so the final
statement follows with their Koszul signs.
\end{proof}
